%% kapa: exam=1 round=23 year=2012 period=1 form=A subject=civil_law subject_ko=민법 seq=33 number=33
%% source: Q-Net 감정평가사 기출문제 (cst00309) artlSeq=5052578 file=raw/1008/5052578/2012년 제23회 감정평가사 제1차시험 기출문제(1교시 A형).pdf
%% conversion: kordoc 4.24.51 -> tools/gen1.py (auto); review=proofread-sonnet; flags=none
\begin{question}{33}
분묘기지권에 관한 설명으로 옳지 않은 것은? (다툼이 있으면 판례에 의함)
\begin{choices}
  \item 분묘기지권은 분묘를 소유하기 위하여 타인의 토지를 제한된 범위 내에서 사용할 수 있는 물권이다.
  \item 분묘가 멸실되더라도 유골이 존재하여 일시적인 멸실에 불과하다면 분묘기지권은 소멸하지 않는다.
  \item 암장(暗葬)되어 있어 객관적으로 인식할 수 있는 외형을 갖추고 있지 아니한 경우에는 분묘기지권이 인정되지 않는다.
  \item 분묘기지권을 시효로 취득한 경우에는 지료를 지급할 필요가 없다.
  \item 분묘기지권의 효력이 미치는 지역의 범위에는 기존의 분묘 외에 새로운 분묘를 신설할 권능도 포함된다.
\end{choices}
\end{question}
\answer{5}
